\SourceLeaf{20}

Die Umkehrung der vorstehenden Entwicklung lehrt, dass man aus den Angaben von \textsc{Gauss} im
Fall $q=3n+1$ zunächst auf die Proportion
\[
\left[\dfrac{\varkappa_1}{\pi_1}\right]:\left[\dfrac{\pi_1}{\varkappa_1}\right]=\left[\dfrac{\varkappa_1}{\pi_2}\right]:\left[\dfrac{\pi_2}{\varkappa_1}\right]
\]
schließen kann. Unter nochmaliger Benutzung des in dieser Proportion ausgesprochenen Gesetzes findet man
\[
\dfrac{\left[\dfrac{\varkappa_1}{\pi_1}\right]}{\left[\dfrac{\pi_1}{\varkappa_1}\right]}=\dfrac{\left[\dfrac{\varkappa_2}{\pi_1}\right]}{\left[\dfrac{\pi_1}{\varkappa_2}\right]}=\dfrac{\left[\dfrac{\varkappa_1}{\pi_2}\right]^2}{\left[\dfrac{\pi_2}{\varkappa_1}\right]^2}
\quad\text{und also}\quad
\dfrac{\left[\dfrac{\varkappa_1}{\pi_2}\right]}{\left[\dfrac{\pi_2}{\varkappa_1}\right]}=\dfrac{\left[\dfrac{\varkappa_1}{\pi_2}\right]^3}{\left[\dfrac{\pi_2}{\varkappa_1}\right]^2}.
\]
Da hieraus $\left[\dfrac{\varkappa_1}{\pi_2}\right]=\left[\dfrac{\pi_2}{\varkappa_1}\right]$ hervorgeht, ist auch für $q=3n+1$ in den Angaben von \textsc{Gauss} das allgemeine Reziprozitätsgesetz in seiner einfachsten Form implizit enthalten.

Unter [4] sind diejenigen Aufzeichnungen gesammelt, durch die \textsc{Gauss} seine kubischen Reziprozitätstheoreme bewiesen hat. Da diese Aufzeichnungen mehreren losen Zetteln ohne Datumsangabe entnommen sind, lässt sich die Entstehungszeit der fraglichen Beweise nicht mit Sicherheit angeben. Es ist
aber wahrscheinlich, dass die Entwicklungen bald nach Gewinnung der »Observatio« und der »Criterien«
durchgeführt wurden. Sie schließen sich nämlich im Gedankengang sehr eng an die ersten Angaben unter [2]
und im Übrigen an Art. 358 der Disquisitiones arithmeticae an. Andererseits war \textsc{Gauss} zufolge seiner
eigenen Aussage\HistoricalNote{\textsuperscript{*)}}{*)}{Siehe Art. 33 der am 24. August 1808 der Göttinger Societät vorgelegten Abhandlung »Summatio quarumdam serierum singularium«.} jedenfalls vor Mitte des Jahres 1808 im Besitz des sechsten Beweises des quadratischen
Reziprozitätsgesetzes, der mit den hier gemeinten Beweisen am nächsten verwandt ist.

Dem Beweis von \textsc{Gauss} für das kubische Reziprozitätsgesetz liegt die Kreisteilung zugrunde. Die
einzelnen Formeln unter [4] können entweder unmittelbar aus Art. 358 der Disquisitiones arithmeticae entnommen werden oder lassen sich zumindest mit sehr geringer Mühe aus den dortigen Entwicklungen ableiten.
Die Kongruenz
\[
P^{pp}\equiv r^{pp}+Rr^{ppg}+\ldots\equiv P(R^{\mathrm{ind.}\,p})
\]
bezieht sich auf den Modul $p$ und ergibt sich daraus, dass alle von 1 verschiedenen Polynomkoeffizienten der Potenz $pp$ durch $p$ teilbar sind. Die Fallunterscheidung
\[
\left\{\dfrac{M-N\sqrt{-27}}{M+N\sqrt{-27}}\right\}^{\frac13(pp-1)}\equiv 1,\quad -\tfrac12+\tfrac12\sqrt{-3},\quad -\tfrac12-\tfrac12\sqrt{-3}\qquad\;(\mathrm{mod.}\ p)
\]
entspricht sowohl für $p=3m-1$ als auch für $p=3m+1$ genau der oben unter I, II und III getroffenen
Unterscheidung $B\equiv 0\;(\mathrm{mod.}\ q)$ usw. Die Zuordnung jener drei Kongruenzen zu den folgenden
\[
p^{\frac13(n-1)}\equiv 1,\quad -\tfrac12-\tfrac12\dfrac{M}{3N},\quad -\tfrac12+\tfrac12\dfrac{M}{3N}\qquad\;(\mathrm{mod.}\ n),
\]
wie sie in [4] entwickelt wird, enthält demnach bereits einen vollständigen Beweis der Angaben von \textsc{Gauss}
und damit implizit des allgemeinen kubischen Reziprozitätsgesetzes.

Die weiteren Ausführungen unter [4] betreffen die Kriterien dafür, dass 2 und 3 kubische Reste sind.
Zum kubischen Rest 3 bemerke man nur, dass (unter Verwendung der in [4] gebrauchten Bezeichnung)
$p^3\equiv p, p'$ oder $p''\;(\mathrm{mod.}\ 3)$ zutrifft, je nachdem, ob ind.\,$3\equiv 0, 1$ oder $2\;(\mathrm{mod.}\ 3)$ gilt. Die Bedingung
$N\equiv 0\;(\mathrm{mod.}\ 3)$ dafür, dass 3 kubischer Rest modulo $n$ ist, ergibt sich daraufhin unmittelbar aus den Schlussangaben von [4].

